% Part: normal-modal-logic
% Chapter: tableaux
% Section: simple-S5

\documentclass[../../../include/open-logic-section]{subfiles}

\begin{document}

\olfileid{nml}{tab}{s5}

\olsection{\Log{S5}-এর সরল \usetoken{P}{tableau}}

\Log{S5} সার্বিক মডেলগুলির শ্রেণি সাপেক্ষে বিশুদ্ধ ও পূর্ণ; অর্থাৎ
সেই মডেলগুলিতে প্রতিটি জগৎ থেকে প্রতিটি জগৎ অভিগম্য। এমন মডেলে
অভিগম্যতা-সম্বন্ধ আলাদা করে বিবেচনা করতে হয় না: ‘এমন একটি জগৎ~$w$
আছে যেখানে $\mSat{M}{!A}[w]$’ তখনই সত্য, যখন~$u$ থেকে অভিগম্য
এমন একটি~$w$ আছে। তাই \Log{S5}-এ মডেলকে শুধু জগৎসমষ্টি ও
মাননির্ধারণ~$V$ দিয়ে সংজ্ঞায়িত করা যায়। এতে ইঙ্গিত মেলে যে
!!{tableau}-র বিধিগুলিও সরল করা যাবে। সাধারণ ক্ষেত্রে
পূর্বসূচি হিসেবে ধনাত্মক পূর্ণসংখ্যার অনুক্রম নিই, যাতে কোন
পূর্বসূচির নামাঙ্কিত জগৎ কোনটি থেকে অভিগম্য, তা জানা যায়:
$\sigma.n$, $\sigma$ থেকে অভিগম্য জগতের নাম দেয়। কিন্তু
\Log{S5}-এ যেকোনো জগৎ অন্য যেকোনো জগৎ থেকে অভিগম্য; তাই
এই হিসাব রাখার দরকার নেই। তার বদলে ধনাত্মক পূর্ণসংখ্যাকেই
পূর্বসূচি হিসেবে ব্যবহার করি। সরলীকৃত বিধিগুলি
\olref{tab:rules-S5}-এ দেওয়া হয়েছে।

\begin{table}
  \begin{center}
    \def\arraystretch{3}\def\fCenter{}
    \begin{tabular}{|c|c|}
    \hline
    \iftag{prvBox}{
      \AxiomC{\sFmla{\True}{\Box !A}[n]}
      \RightLabel{\TRule{\True}{\Box}}
      \UnaryInfC{\sFmla{\True}{!A}[m]}
      \DisplayProof
      &
      \AxiomC{\sFmla{\False}{\Box !A}[n]}
      \RightLabel{\TRule{\False}{\Box}}
      \UnaryInfC{\sFmla{\False}{!A}[m]}
      \DisplayProof\\
      $m$ ব্যবহৃত & $m$ নতুন
      \\[1ex]
      \hline}{}
    \iftag{prvDiamond}{
      \AxiomC{\sFmla{\True}{\Diamond !A}[n]}
      \RightLabel{\TRule{\True}{\Diamond}}
      \UnaryInfC{\sFmla{\True}{!A}[m]}
      \DisplayProof
      &
      \AxiomC{\sFmla{\False}{\Diamond !A}[n]}
      \RightLabel{\TRule{\False}{\Diamond}}
      \UnaryInfC{\sFmla{\False}{!A}[m]}
      \DisplayProof\\
      $m$ নতুন & $m$ ব্যবহৃত
      \\[1ex]
      \hline}{}
    \end{tabular}
  \end{center}
  \caption{\Log{S5}-এর সরলীকৃত বিধি।}
  \ollabel{tab:rules-S5}
\end{table}

\begin{ex}
  $\Log{S5} \Proves \Ax{5}$, অর্থাৎ
  $\Diamond!A \lif \Box\Diamond!A$, দেখায় এমন একটি
  সরলীকৃত বদ্ধ ট্যাবলো দিচ্ছি।
  \begin{oltableau}
    [\pFmla{\False}{\Diamond\formula{A} \lif \Box\Diamond \formula{A}}{1},
      just = \TAss
      [\pFmla{\True}{\Diamond \formula{A}}{1}, just = {\TRule{\False}{\lif}[1]}
        [\pFmla{\False}{\Box\Diamond \formula{A}}{1},
          just = {\TRule{\False}{\lif}[1]}
          [\pFmla{\False}{\Diamond \formula{A}}{2},
            just = {\TRule{\False}{\Box}[3]}
            [\pFmla{\True}{\formula{A}}{3},
              just = {\TRule{\True}{\Diamond}[2]}
                [\pFmla{\False}{\formula{A}}{3},
                  just = {\TRule{\False}{\Diamond}[4]}, close]
            ]
          ]
        ]
      ]
    ]
  \end{oltableau}
\end{ex}

\end{document}
